% =============================================================
% CHAPTER 3 — RELATED WORK / STATE OF THE ART
% Purpose: show you know the field, organise prior work into themes, and -- most
% importantly -- identify the GAP your thesis fills. This is not a list of
% summaries; it is an argument that leads to "...and that is why this thesis is
% needed."
% Tip: group by theme or approach, NOT one paragraph per paper. Compare and
% contrast. Every subsection should end with a sentence about limitations/gaps.
% =============================================================
\chapter{Related Work}
\label{chap:relatedwork}

\guide{Open with one paragraph describing how you organised the literature (the
themes below) and what question you are using it to answer: ``what has been tried
for face-swap detection, and where do MLLMs fit?''}

\section{Classical and Deep-Learning Face-Swap Detection}
\guide{Review the main line of dedicated detectors: CNN-based classifiers,
frequency/artefact-based methods, and benchmarks (e.g. FaceForensics++,
Celeb-DF). Characterise their strengths (high in-domain accuracy) and the
recurring weakness you care about: poor generalisation to unseen manipulation
methods. This sets up why a more general detector is attractive.}

\section{Generalisation and Cross-Dataset Detection}
\guide{Discuss work specifically about the generalisation problem -- detectors
that fail on manipulation types or datasets they were not trained on. This is the
pain point that motivates zero-shot MLLM use.}

\section{(M)LLMs for Vision Understanding and Forensics}
\guide{Review the emerging work that uses vision-language models / MLLMs for
image understanding, reasoning, and especially for detecting manipulated or
AI-generated images. Be precise about what each did and how your study differs
(e.g. different datasets, broader model coverage, systematic prompt study,
reproducible pipeline).}

\section{Datasets and Benchmarks}
\guide{Introduce the datasets used in the literature and position the ones you
use (FOWS, GOTCHA) among them: what they contain, why they are relevant, and any
known limitations. (Full details of how YOU use them go in \Cref{chap:setup}.)}

\section{Summary and Research Gap}
\guide{Close the chapter by stating the gap in one tight paragraph: what nobody
has adequately done yet, and how your research questions (\Cref{chap:introduction})
target exactly that gap. This paragraph is the bridge into your methodology.}
